

% 02 Apr: Clean up

\section{Visible Congregation}

Editorial article in “Folkebladet” for January 20, 1897.

In Free Church work the congregation is of such surpassing importance that we cannot labor enough for it to be acknowledged in its truth according to God’s Word.

In the New Testament such high and glorious words are spoken concerning the congregation that many have given it up in despair and have sought to escape the holy demand by teaching that all the great and glorious things are said of an invisible congregation. But what we ordinarily have under the name of congregation, that is only the visible congregation, and that is quite another matter.

But what if according to the New Testament it were the same matter?

It is not uncommon in our cities to see a number of members of the congregation—belonging to Norwegian Lutheran congregations—go drunken, reeling with brandy, from saloon to saloon; such, then, are taken for the ‘visible’ congregation.

And conversely, when a congregation is altogether sunk in worldliness and indifference, so that one knows of not a single Christian in the congregation, then one consoles oneself with the assurance that there is at least an invisible congregation somewhere, if nowhere else than in the newborn children.

Where does such a thing lead? Invariably to this: that one despises the congregation and exalts the Church at the congregation’s expense. One sees that the local congregation is neither light nor salt; neither believing nor living; it surely cannot possibly be Christ’s body and Christ’s bride, as it is. And then, in order to escape the truth and the honest confession of sin and transgression, one forms for oneself a dim and vague notion that there is, lifted high above the congregation, a Church about which one prefers not to say anything definite—where it is to be found, or where it ever was.

If there is to be a change for the better among us, then we must begin by acknowledging that the Church is the congregations; and that as the congregations are, so is the Church, neither better nor worse, for it is precisely one and the same matter. We must cease deceiving ourselves with the notion that if we have a bad congregation, we nevertheless have a good and glorious Church; for the Church is altogether nothing other than the congregations, and a bad congregation is shame and misery for itself and its Lord and Savior, and it is harmful to other congregations, whether few or many, according as it comes into contact with them.

But next, it is above all important that it be acknowledged among us that our congregations do not become visible through unbelieving, immoral, profane, and drink-sodden church members; such serve only to hide, cover over, and disfigure the congregation, just as a fair face is disfigured when it is covered all over with filthy splashes.

It is of the very highest importance that the congregation be visible. For it is Jesus Christ’s own body, set to reveal him in the world. What is invisible cannot make manifest; light is required. Therefore what matters is that it may be said of the congregation: “You were formerly darkness, but now you are light in the Lord.” If that cannot be said, then it is no visible congregation.

In other words, the congregation becomes visible by its believing members and their activity.

It becomes visible by God’s Word, that is true, for the Word is a light for the foot and a lantern on the path; its proclamation and use in church and house, in testimony and conversation, serve to make the congregation visible.

It becomes visible through Baptism and the Lord’s Supper; for in Baptism one is clothed with Christ, and in the Supper his death is proclaimed until he comes.

But it is not merely a matter of the Word and the Sacrament without faith and without believing people in the congregation. It is written: You are the light of the world, you are the salt of the earth. Scripture does not merely say that God’s Word is the light of the world.

Christ’s body—that through which Christ is seen, that through which Christ works—is visible, because it is believing men and women who labor with the Word and the Sacraments, and it is their outward, organized community which is called the congregation.

But that does not fit our outwardly organized congregations at all! No, but from that it only follows that we must ask: Do you want a congregation? If the congregation has been lost to us in the undifferentiated mass of the state church, so that we call that congregation which is darkness and not light, then there is simply nothing else to be done but to obey the command: Awake, you who sleep, arise from the dead, and Christ shall shine upon you!

This is the Free Church’s first task: to bring the congregation out of the rubble, or, as the Prophet says, to rebuild ruins and restore desolate places.

It is not the great body that matters now, but each individual little congregation. The matter is that in every place where baptized people gather around the Word, there comes to be a visible congregation, a city on a hill, which cannot be hidden.

Men and women, let us labor with faithfulness toward this goal, each in our own place and all together, as many as have the faith of Jesus and the life in the Son of God.

Georg Sverdrup
